% Complément expérimental de la soumission EGC 2027 (version anonyme).
% Compilation : pdflatex complements.tex (deux fois pour les renvois).
\documentclass[a4paper,10pt]{article}
\usepackage[T1]{fontenc}
\usepackage[utf8]{inputenc}
\usepackage[french]{babel}
\usepackage{lmodern}
\usepackage[margin=2.2cm]{geometry}
\usepackage{amsmath}
\usepackage{amssymb}
\usepackage{booktabs}
\usepackage{icomma}
\usepackage[hidelinks,pdfauthor={},pdftitle={Complément expérimental}]{hyperref}
\renewcommand{\thetable}{S\arabic{table}}
\setcounter{secnumdepth}{1}

\title{Complément expérimental\\[0.4em]\large Apprentissage de projections de correction : spécialiser a posteriori un espace d'embedding sur une base de données}
\author{Soumission anonyme}
\date{}

\begin{document}
\maketitle

\section{Contenu et périmètre}

Ce complément rassemble les résultats détaillés qui ne tiennent pas dans l'article. Les expériences portent sur quatre jeux MTEB (SciFact, NFCorpus, ArguAna et SciDocs) et cinq encodeurs, soit vingt couples jeu de données--encodeur. Un \emph{découpage} désigne la répartition des requêtes en ensembles d'entraînement, de validation et de test (60/20/20), stratifiée selon leur statut (réussie ou en erreur).

Sur le découpage initial, le protocole comprend 520 exécutions : 480 exécutions de correction, dont 468 ont abouti, et 40 exécutions des deux références non supervisées. Les 12 exécutions manquantes, par clustering sur Gemini--SciFact, ne peuvent pas être exécutées, car le nombre d'erreurs résiduelles est inférieur au nombre de clusters demandé.

Les tableaux ne portent pas tous sur les mêmes requêtes :
\begin{itemize}
  \item le tableau~\ref{tab:s-validation-croisee} présente les résultats de test sur cinq découpages, selon le protocole du tableau~1 de l'article ;
  \item les tableaux~\ref{tab:s-references-synthese} à \ref{tab:s-references-couples-2} et \ref{tab:s-covariables} portent sur le test du découpage initial ;
  \item les tableaux d'ablation~\ref{tab:s-ablation-arguana} à \ref{tab:s-clusters-scifact} présentent des variations calculées sur l'ensemble des requêtes du découpage initial, entraînement, validation et test confondus. Ils décrivent l'effet des réglages sur les configurations obtenues, mais ne mesurent pas la généralisation : ils ne doivent pas être comparés directement aux résultats de test.
\end{itemize}

Les notations suivent l'article. Var. : variante globale (Q) ou par clustering (C) ; \textit{gate} : critère d'acceptation strict (S) ou adaptatif (A) ; SS : \textit{search strictness}, c'est-à-dire le poids de préservation $w_h$, au niveau \textit{error-scaled} (ES), \textit{mid} (M) ou \textit{relaxed} (R). Dég. compte les requêtes initialement réussies devenues en erreur ; Fix@1 compte les requêtes initialement en erreur corrigées en top-1.

\section{Règle de sélection des configurations}

Chaque couple compte 24 exécutions de correction abouties (12 sur Gemini--SciFact). Celle présentée pour le découpage initial est obtenue en comparant, de gauche à droite, les critères suivants calculés en validation, en faveur de la valeur la plus élevée :
\begin{equation*}
  \bigl(\mathbb{1}[F_{val}\geq D_{val}],\ \Delta\text{NDCG}_{val},\ \Delta\text{MRR}_{val},\ G_{val},\ F_{val},\ T_{val}\bigr),
\end{equation*}
où $D_{val}$ compte les requêtes réussies dégradées, $F_{val}$ les requêtes en erreur corrigées en top-1, $G_{val}$ le gain de rang total et $T_{val}$ les entrées dans le top-10. Le premier terme est la contrainte de \textit{net correction} : une exécution doit corriger en top-1 au moins autant de requêtes en erreur qu'elle dégrade de requêtes réussies. L'inégalité large laisse éligibles les exécutions sans dégradation qui ne corrigent aucune erreur en top-1 ; leurs gains proviennent alors d'entrées dans le top-10. Treize des vingt configurations retenues n'ont aucune dégradation en validation ; les sept autres en acceptent un nombre limité en échange de corrections plus nombreuses. Le test n'intervient jamais dans cette sélection.

Cette règle compare des exécutions terminées. Elle se distingue du critère d'acceptation décrit dans l'article, qui compare les états intermédiaires (\textit{checkpoints}) au sein d'une exécution et privilégie le gain de rang plutôt que le NDCG@10.

\section{Validation croisée par découpages répétés}

Pour chaque couple, les trois meilleures configurations distinctes selon la règle précédente sont réexécutées sur quatre nouveaux découpages aléatoires, soit 240 exécutions supplémentaires, dont 238 ont abouti. Les deux exécutions par clustering manquantes sur SciFact--Gemini s'arrêtent, comme prévu par l'implémentation, lorsque le nouveau découpage d'entraînement laisse un cluster d'erreurs réduit à une seule requête. Sur chaque nouveau découpage, la même règle sélectionne une configuration en validation, puis évalue ce choix en test. Le tableau~\ref{tab:s-validation-croisee} agrège ces quatre résultats et celui du découpage initial.

Un test unilatéral des rangs signés de Wilcoxon sur les vingt gains moyens de NDCG@10 par couple confirme la prédominance des gains positifs ($W=198$, $p=6{,}7\times10^{-5}$). Plusieurs couples partageant un jeu de données ou un encodeur, l'indépendance des observations n'est toutefois pas assurée.

\begin{table}[p]
\centering
\small
\begin{tabular}{lllccll}
\toprule
Jeu & Encodeur & Découpages & $\Delta$NDCG@10 test & $\Delta$MRR test & Pos & Cfg \\
\midrule
ArguAna & Titan-v1 & 5 & +0,0095 $\pm$ 0,0033 & +0,0059 $\pm$ 0,0020 & 5/5 & 1 \\
ArguAna & Gemini-001 & 5 & +0,1224 $\pm$ 0,0284 & +0,1496 $\pm$ 0,0390 & 5/5 & 2 \\
ArguAna & OpenAI-3-large & 5 & +0,0509 $\pm$ 0,0043 & +0,0403 $\pm$ 0,0025 & 5/5 & 3 \\
ArguAna & OpenAI-3-small & 5 & +0,0424 $\pm$ 0,0087 & +0,0355 $\pm$ 0,0033 & 5/5 & 3 \\
ArguAna & MiniLM-L6 & 5 & +0,0155 $\pm$ 0,0036 & +0,0140 $\pm$ 0,0021 & 5/5 & 3 \\
\midrule
NFCorpus & Titan-v1 & 5 & +0,0026 $\pm$ 0,0033 & $-$0,0012 $\pm$ 0,0031 & 4/5 & 2 \\
NFCorpus & Gemini-001 & 5 & +0,0004 $\pm$ 0,0037 & +0,0007 $\pm$ 0,0033 & 3/5 & 3 \\
NFCorpus & OpenAI-3-large & 5 & +0,0009 $\pm$ 0,0017 & +0,0007 $\pm$ 0,0033 & 3/5 & 2 \\
NFCorpus & OpenAI-3-small & 5 & $-$0,0009 $\pm$ 0,0025 & +0,0029 $\pm$ 0,0019 & 2/5 & 2 \\
NFCorpus & MiniLM-L6 & 5 & +0,0048 $\pm$ 0,0044 & +0,0031 $\pm$ 0,0036 & 5/5 & 2 \\
\midrule
SciDocs & Titan-v1 & 5 & +0,0008 $\pm$ 0,0026 & +0,0045 $\pm$ 0,0046 & 4/5 & 3 \\
SciDocs & Gemini-001 & 5 & +0,0171 $\pm$ 0,0062 & +0,0219 $\pm$ 0,0083 & 5/5 & 3 \\
SciDocs & OpenAI-3-large & 5 & +0,0055 $\pm$ 0,0032 & +0,0101 $\pm$ 0,0073 & 5/5 & 2 \\
SciDocs & OpenAI-3-small & 5 & +0,0041 $\pm$ 0,0037 & +0,0061 $\pm$ 0,0070 & 4/5 & 3 \\
SciDocs & MiniLM-L6 & 5 & $-$0,0027 $\pm$ 0,0013 & $-$0,0018 $\pm$ 0,0045 & 0/5 & 3 \\
\midrule
SciFact & Titan-v1 & 5 & +0,0147 $\pm$ 0,0146 & +0,0129 $\pm$ 0,0133 & 5/5 & 2 \\
SciFact & Gemini-001 & 5 & $-$0,0003 $\pm$ 0,0011 & $-$0,0004 $\pm$ 0,0016 & 1/5 & 3 \\
SciFact & OpenAI-3-large & 5 & +0,0194 $\pm$ 0,0051 & +0,0214 $\pm$ 0,0036 & 5/5 & 3 \\
SciFact & OpenAI-3-small & 5 & +0,0033 $\pm$ 0,0022 & +0,0038 $\pm$ 0,0023 & 5/5 & 2 \\
SciFact & MiniLM-L6 & 5 & +0,0108 $\pm$ 0,0085 & +0,0109 $\pm$ 0,0083 & 5/5 & 2 \\
\bottomrule
\end{tabular}
\caption{Validation croisée par découpages aléatoires répétés. Pour chaque couple, les trois meilleures configurations en validation sur le découpage initial sont réexécutées sur quatre nouveaux découpages ; sur chacun, la configuration retenue en validation est évaluée en test. Les valeurs agrègent ces quatre résultats et celui du découpage initial (moyenne $\pm$ écart-type). Pos : découpages avec un gain de NDCG@10 en test ; Cfg : configurations différentes retenues.}
\label{tab:s-validation-croisee}
\end{table}


\section{Comparaison aux références non supervisées}

Les références \textit{whitening} et \textit{all-but-the-top} sont ajustées sur les embeddings du corpus et appliquées au corpus seul, avec les mêmes découpages et métriques que la méthode. Le tableau~\ref{tab:s-references-synthese} résume la comparaison par jeu ; les tableaux~\ref{tab:s-references-couples-1} et \ref{tab:s-references-couples-2} détaillent les résultats par couple. Les lignes « correction (sélection) » correspondent à la configuration retenue sur le découpage initial.

\begin{table}[p]
\centering
\small
\begin{tabular}{llrrrrr}
\toprule
Jeu & Méthode & $\Delta$NDCG@10 & $\Delta$MRR & Couples & Dég. & Fix@1 \\
\midrule
ArguAna & correction (sélection) & +0,0579 & +0,0614 & 5/5 & 0 & 93 \\
ArguAna & \textit{all-but-the-top} & $-$0,0125 & $-$0,0119 & 1/5 & 22 & 0 \\
ArguAna & \textit{whitening} & $-$0,0162 & $-$0,0179 & 2/5 & 27 & 0 \\
\midrule
NFCorpus & correction (sélection) & +0,0015 & +0,0003 & 4/5 & 34 & 34 \\
NFCorpus & \textit{all-but-the-top} & $-$0,0202 & $-$0,0218 & 1/5 & 154 & 92 \\
NFCorpus & \textit{whitening} & $-$0,1649 & $-$0,1605 & 0/5 & 508 & 120 \\
\midrule
SciDocs & correction (sélection) & +0,0068 & +0,0095 & 4/5 & 11 & 20 \\
SciDocs & \textit{all-but-the-top} & $-$0,0074 & $-$0,0062 & 1/5 & 32 & 32 \\
SciDocs & \textit{whitening} & $-$0,0412 & $-$0,0564 & 0/5 & 81 & 42 \\
\midrule
SciFact & correction (sélection) & +0,0149 & +0,0142 & 4/5 & 4 & 17 \\
SciFact & \textit{all-but-the-top} & $-$0,0136 & $-$0,0142 & 0/5 & 27 & 15 \\
SciFact & \textit{whitening} & $-$0,0591 & $-$0,0521 & 0/5 & 63 & 29 \\
\bottomrule
\end{tabular}
\caption{Synthèse par jeu de la comparaison aux références non supervisées, sur le test du découpage initial. $\Delta$NDCG@10 et $\Delta$MRR : moyennes sur les cinq encodeurs ; Couples : encodeurs avec un gain de NDCG@10 en test ; Dég. et Fix@1 : totaux sur les cinq encodeurs des requêtes initialement réussies devenues en erreur et des requêtes initialement en erreur corrigées en top-1.}
\label{tab:s-references-synthese}
\end{table}

\begin{table}[p]
\centering
\footnotesize
\begin{tabular}{lllcrrrr}
\toprule
Jeu & Encodeur & Méthode & Base & $\Delta$NDCG@10 & $\Delta$MRR & Dég. & Fix@1 \\
\midrule
ArguAna & Titan-v1 & correction (sélection) & 0,37/0,26 & +0,0116 & +0,0081 & 0 & 0 \\
ArguAna & Titan-v1 & \textit{all-but-the-top} & 0,37/0,26 & $-$0,0033 & +0,0001 & 0 & 0 \\
ArguAna & Titan-v1 & \textit{whitening} & 0,37/0,26 & $-$0,0072 & $-$0,0075 & 0 & 0 \\
ArguAna & Gemini-001 & correction (sélection) & 0,58/0,46 & +0,1626 & +0,2063 & 0 & 93 \\
ArguAna & Gemini-001 & \textit{all-but-the-top} & 0,58/0,46 & $-$0,0663 & $-$0,0685 & 22 & 0 \\
ArguAna & Gemini-001 & \textit{whitening} & 0,58/0,46 & $-$0,0740 & $-$0,0829 & 27 & 0 \\
ArguAna & OpenAI-3-large & correction (sélection) & 0,41/0,28 & +0,0535 & +0,0401 & 0 & 0 \\
ArguAna & OpenAI-3-large & \textit{all-but-the-top} & 0,41/0,28 & +0,0000 & +0,0041 & 0 & 0 \\
ArguAna & OpenAI-3-large & \textit{whitening} & 0,41/0,28 & +0,0038 & $-$0,0023 & 0 & 0 \\
ArguAna & OpenAI-3-small & correction (sélection) & 0,40/0,28 & +0,0506 & +0,0401 & 0 & 0 \\
ArguAna & OpenAI-3-small & \textit{all-but-the-top} & 0,40/0,28 & +0,0094 & +0,0081 & 0 & 0 \\
ArguAna & OpenAI-3-small & \textit{whitening} & 0,40/0,28 & +0,0038 & +0,0068 & 0 & 0 \\
ArguAna & MiniLM-L6 & correction (sélection) & 0,37/0,26 & +0,0113 & +0,0126 & 0 & 0 \\
ArguAna & MiniLM-L6 & \textit{all-but-the-top} & 0,37/0,26 & $-$0,0025 & $-$0,0031 & 0 & 0 \\
ArguAna & MiniLM-L6 & \textit{whitening} & 0,37/0,26 & $-$0,0072 & $-$0,0035 & 0 & 0 \\
\midrule
NFCorpus & Titan-v1 & correction (sélection) & 0,33/0,54 & +0,0043 & $-$0,0027 & 6 & 2 \\
NFCorpus & Titan-v1 & \textit{all-but-the-top} & 0,33/0,54 & +0,0078 & +0,0042 & 10 & 12 \\
NFCorpus & Titan-v1 & \textit{whitening} & 0,33/0,54 & $-$0,0399 & $-$0,0168 & 27 & 25 \\
NFCorpus & Gemini-001 & correction (sélection) & 0,41/0,62 & +0,0041 & +0,0039 & 7 & 9 \\
NFCorpus & Gemini-001 & \textit{all-but-the-top} & 0,41/0,62 & $-$0,0377 & $-$0,0344 & 37 & 20 \\
NFCorpus & Gemini-001 & \textit{whitening} & 0,41/0,62 & $-$0,2403 & $-$0,2255 & 140 & 22 \\
NFCorpus & OpenAI-3-large & correction (sélection) & 0,40/0,61 & +0,0012 & $-$0,0008 & 6 & 6 \\
NFCorpus & OpenAI-3-large & \textit{all-but-the-top} & 0,40/0,61 & $-$0,0379 & $-$0,0372 & 46 & 23 \\
NFCorpus & OpenAI-3-large & \textit{whitening} & 0,40/0,61 & $-$0,2980 & $-$0,3612 & 188 & 13 \\
NFCorpus & OpenAI-3-small & correction (sélection) & 0,36/0,57 & $-$0,0037 & +0,0010 & 11 & 11 \\
NFCorpus & OpenAI-3-small & \textit{all-but-the-top} & 0,36/0,57 & $-$0,0270 & $-$0,0240 & 36 & 21 \\
NFCorpus & OpenAI-3-small & \textit{whitening} & 0,36/0,57 & $-$0,1816 & $-$0,1392 & 88 & 27 \\
NFCorpus & MiniLM-L6 & correction (sélection) & 0,32/0,51 & +0,0017 & +0,0000 & 4 & 6 \\
NFCorpus & MiniLM-L6 & \textit{all-but-the-top} & 0,32/0,51 & $-$0,0064 & $-$0,0178 & 25 & 16 \\
NFCorpus & MiniLM-L6 & \textit{whitening} & 0,32/0,51 & $-$0,0647 & $-$0,0597 & 65 & 33 \\
\bottomrule
\end{tabular}
\caption{Comparaison par couple aux références non supervisées, sur le test du découpage initial (partie 1/2). Les références appliquent une transformation globale unique au corpus, sans supervision ni critère d'acceptation. Base : NDCG@10/MRR avant correction ; Dég. et Fix@1 : requêtes initialement réussies devenues en erreur et requêtes initialement en erreur corrigées en top-1.}
\label{tab:s-references-couples-1}
\end{table}

\begin{table}[p]
\centering
\footnotesize
\begin{tabular}{lllcrrrr}
\toprule
Jeu & Encodeur & Méthode & Base & $\Delta$NDCG@10 & $\Delta$MRR & Dég. & Fix@1 \\
\midrule
SciDocs & Titan-v1 & correction (sélection) & 0,21/0,36 & +0,0042 & +0,0095 & 0 & 3 \\
SciDocs & Titan-v1 & \textit{all-but-the-top} & 0,21/0,36 & +0,0018 & +0,0073 & 1 & 3 \\
SciDocs & Titan-v1 & \textit{whitening} & 0,21/0,36 & $-$0,0051 & $-$0,0188 & 6 & 2 \\
SciDocs & Gemini-001 & correction (sélection) & 0,24/0,40 & +0,0220 & +0,0190 & 4 & 7 \\
SciDocs & Gemini-001 & \textit{all-but-the-top} & 0,24/0,40 & $-$0,0085 & +0,0008 & 8 & 10 \\
SciDocs & Gemini-001 & \textit{whitening} & 0,24/0,40 & $-$0,0536 & $-$0,0570 & 17 & 13 \\
SciDocs & OpenAI-3-large & correction (sélection) & 0,22/0,38 & +0,0099 & +0,0156 & 4 & 7 \\
SciDocs & OpenAI-3-large & \textit{all-but-the-top} & 0,22/0,38 & $-$0,0090 & $-$0,0128 & 9 & 11 \\
SciDocs & OpenAI-3-large & \textit{whitening} & 0,22/0,38 & $-$0,0661 & $-$0,1069 & 24 & 9 \\
SciDocs & OpenAI-3-small & correction (sélection) & 0,19/0,34 & +0,0018 & $-$0,0009 & 0 & 0 \\
SciDocs & OpenAI-3-small & \textit{all-but-the-top} & 0,19/0,34 & $-$0,0088 & $-$0,0189 & 9 & 6 \\
SciDocs & OpenAI-3-small & \textit{whitening} & 0,19/0,34 & $-$0,0520 & $-$0,0796 & 24 & 9 \\
SciDocs & MiniLM-L6 & correction (sélection) & 0,20/0,35 & $-$0,0040 & +0,0044 & 3 & 3 \\
SciDocs & MiniLM-L6 & \textit{all-but-the-top} & 0,20/0,35 & $-$0,0123 & $-$0,0072 & 5 & 2 \\
SciDocs & MiniLM-L6 & \textit{whitening} & 0,20/0,35 & $-$0,0293 & $-$0,0198 & 10 & 9 \\
\midrule
SciFact & Titan-v1 & correction (sélection) & 0,73/0,70 & +0,0323 & +0,0279 & 1 & 5 \\
SciFact & Titan-v1 & \textit{all-but-the-top} & 0,73/0,70 & $-$0,0079 & +0,0029 & 1 & 4 \\
SciFact & Titan-v1 & \textit{whitening} & 0,73/0,70 & $-$0,0301 & $-$0,0162 & 3 & 4 \\
SciFact & Gemini-001 & correction (sélection) & 0,99/0,99 & +0,0000 & +0,0000 & 0 & 0 \\
SciFact & Gemini-001 & \textit{all-but-the-top} & 0,99/0,99 & $-$0,0023 & $-$0,0031 & 1 & 0 \\
SciFact & Gemini-001 & \textit{whitening} & 0,99/0,99 & $-$0,0446 & $-$0,0412 & 10 & 0 \\
SciFact & OpenAI-3-large & correction (sélection) & 0,77/0,74 & +0,0278 & +0,0272 & 0 & 6 \\
SciFact & OpenAI-3-large & \textit{all-but-the-top} & 0,77/0,74 & $-$0,0189 & $-$0,0269 & 10 & 4 \\
SciFact & OpenAI-3-large & \textit{whitening} & 0,77/0,74 & $-$0,1191 & $-$0,1174 & 22 & 7 \\
SciFact & OpenAI-3-small & correction (sélection) & 0,72/0,69 & +0,0065 & +0,0028 & 2 & 2 \\
SciFact & OpenAI-3-small & \textit{all-but-the-top} & 0,72/0,69 & $-$0,0176 & $-$0,0280 & 10 & 4 \\
SciFact & OpenAI-3-small & \textit{whitening} & 0,72/0,69 & $-$0,0766 & $-$0,0714 & 19 & 9 \\
SciFact & MiniLM-L6 & correction (sélection) & 0,65/0,62 & +0,0077 & +0,0130 & 1 & 4 \\
SciFact & MiniLM-L6 & \textit{all-but-the-top} & 0,65/0,62 & $-$0,0211 & $-$0,0160 & 5 & 3 \\
SciFact & MiniLM-L6 & \textit{whitening} & 0,65/0,62 & $-$0,0251 & $-$0,0145 & 9 & 9 \\
\bottomrule
\end{tabular}
\caption{Comparaison par couple aux références non supervisées, sur le test du découpage initial (partie 2/2). Les références appliquent une transformation globale unique au corpus, sans supervision ni critère d'acceptation. Base : NDCG@10/MRR avant correction ; Dég. et Fix@1 : requêtes initialement réussies devenues en erreur et requêtes initialement en erreur corrigées en top-1.}
\label{tab:s-references-couples-2}
\end{table}


\section{Ablations}

Les tableaux~\ref{tab:s-ablation-arguana} à \ref{tab:s-ablation-scifact} croisent les deux \textit{gates} et les trois niveaux de \textit{search strictness}, pour chaque variante et chaque encodeur. Les tableaux~\ref{tab:s-clusters-arguana} à \ref{tab:s-clusters-scifact} présentent l'effet du nombre de clusters de la variante par clustering. Rappel : ces variations sont calculées sur l'ensemble des requêtes du découpage initial.

\begin{table}[p]
\centering
\small
\begin{tabular}{lllrrrrrr}
\toprule
 & & & \multicolumn{3}{c}{\textit{Gate} strict} & \multicolumn{3}{c}{\textit{Gate} adaptatif} \\
\cmidrule(lr){4-6}\cmidrule(lr){7-9}
Encodeur & Var. & SS & $\Delta$NDCG@10 & $\Delta$MRR & Dég. & $\Delta$NDCG@10 & $\Delta$MRR & Dég. \\
\midrule
Titan-v1 & Q & ES & +0,0017 & +0,0012 & 0 & +0,0000 & +0,0000 & 0 \\
Titan-v1 & Q & M & +0,0015 & +0,0010 & 0 & +0,0000 & +0,0000 & 0 \\
Titan-v1 & Q & R & +0,0020 & +0,0013 & 0 & +0,0000 & +0,0000 & 0 \\
Titan-v1 & C & ES & +0,0354 & +0,0204 & 0 & +0,0000 & +0,0000 & 0 \\
Titan-v1 & C & M & +0,0508 & +0,0295 & 0 & +0,0000 & +0,0000 & 0 \\
Titan-v1 & C & R & +0,0522 & +0,0313 & 0 & +0,0000 & +0,0000 & 0 \\
\midrule
Gemini-001 & Q & ES & +0,0213 & +0,0245 & 6 & +0,0907 & +0,1133 & 9 \\
Gemini-001 & Q & M & +0,0426 & +0,0512 & 5 & +0,0851 & +0,1055 & 10 \\
Gemini-001 & Q & R & +0,0387 & +0,0450 & 7 & +0,1150 & +0,1417 & 7 \\
Gemini-001 & C & ES & +0,0000 & +0,0000 & 0 & +0,1459 & +0,1832 & 8 \\
Gemini-001 & C & M & +0,0367 & +0,0441 & 5 & +0,0976 & +0,1207 & 0 \\
Gemini-001 & C & R & +0,0525 & +0,0617 & 6 & +0,1149 & +0,1406 & 9 \\
\midrule
OpenAI-3-large & Q & ES & +0,0172 & +0,0125 & 0 & +0,0000 & +0,0000 & 0 \\
OpenAI-3-large & Q & M & +0,0172 & +0,0128 & 0 & +0,0000 & +0,0000 & 0 \\
OpenAI-3-large & Q & R & +0,0172 & +0,0129 & 0 & +0,0000 & +0,0000 & 0 \\
OpenAI-3-large & C & ES & +0,0736 & +0,0573 & 0 & +0,0000 & +0,0000 & 0 \\
OpenAI-3-large & C & M & +0,0872 & +0,0681 & 0 & +0,0000 & +0,0000 & 0 \\
OpenAI-3-large & C & R & +0,0857 & +0,0677 & 0 & +0,0000 & +0,0000 & 0 \\
\midrule
OpenAI-3-small & Q & ES & +0,0151 & +0,0109 & 0 & +0,0000 & +0,0000 & 0 \\
OpenAI-3-small & Q & M & +0,0152 & +0,0108 & 0 & +0,0000 & +0,0000 & 0 \\
OpenAI-3-small & Q & R & +0,0151 & +0,0107 & 0 & +0,0000 & +0,0000 & 0 \\
OpenAI-3-small & C & ES & +0,0814 & +0,0635 & 0 & +0,0000 & +0,0000 & 0 \\
OpenAI-3-small & C & M & +0,0784 & +0,0631 & 0 & +0,0000 & +0,0000 & 0 \\
OpenAI-3-small & C & R & +0,0820 & +0,0652 & 0 & +0,0000 & +0,0000 & 0 \\
\midrule
MiniLM-L6 & Q & ES & +0,0106 & +0,0067 & 0 & +0,0000 & +0,0000 & 0 \\
MiniLM-L6 & Q & M & +0,0105 & +0,0069 & 0 & +0,0000 & +0,0000 & 0 \\
MiniLM-L6 & Q & R & +0,0106 & +0,0068 & 0 & +0,0000 & +0,0000 & 0 \\
MiniLM-L6 & C & ES & +0,0536 & +0,0366 & 0 & +0,0000 & +0,0000 & 0 \\
MiniLM-L6 & C & M & +0,0482 & +0,0322 & 0 & +0,0000 & +0,0000 & 0 \\
MiniLM-L6 & C & R & +0,0592 & +0,0399 & 0 & +0,0000 & +0,0000 & 0 \\
\bottomrule
\end{tabular}
\caption{Ablation du \textit{gate} et de la \textit{search strictness}, ArguAna. Variations calculées sur l'ensemble des requêtes du découpage initial. Var. : Q globale, C par clustering ; SS : \textit{search strictness} (ES \textit{error-scaled}, M \textit{mid}, R \textit{relaxed}) ; Dég. : requêtes initialement réussies devenues en erreur.}
\label{tab:s-ablation-arguana}
\end{table}

\begin{table}[p]
\centering
\small
\begin{tabular}{lllrrrrrr}
\toprule
 & & & \multicolumn{3}{c}{\textit{Gate} strict} & \multicolumn{3}{c}{\textit{Gate} adaptatif} \\
\cmidrule(lr){4-6}\cmidrule(lr){7-9}
Encodeur & Var. & SS & $\Delta$NDCG@10 & $\Delta$MRR & Dég. & $\Delta$NDCG@10 & $\Delta$MRR & Dég. \\
\midrule
Titan-v1 & Q & ES & +0,0000 & +0,0000 & 0 & +0,0040 & +0,0018 & 14 \\
Titan-v1 & Q & M & +0,0014 & +0,0005 & 17 & +0,0001 & +0,0005 & 10 \\
Titan-v1 & Q & R & +0,0000 & +0,0000 & 0 & +0,0002 & +0,0005 & 8 \\
Titan-v1 & C & ES & +0,0000 & +0,0000 & 0 & +0,0043 & +0,0008 & 19 \\
Titan-v1 & C & M & +0,0018 & +0,0022 & 17 & +0,0001 & +0,0010 & 5 \\
Titan-v1 & C & R & +0,0000 & +0,0000 & 0 & +0,0000 & +0,0000 & 0 \\
\midrule
Gemini-001 & Q & ES & +0,0000 & +0,0000 & 0 & +0,0010 & +0,0021 & 6 \\
Gemini-001 & Q & M & +0,0000 & +0,0000 & 0 & +0,0000 & +0,0000 & 0 \\
Gemini-001 & Q & R & +0,0000 & +0,0011 & 8 & +0,0031 & +0,0046 & 22 \\
Gemini-001 & C & ES & +0,0000 & +0,0000 & 0 & +0,0041 & +0,0046 & 28 \\
Gemini-001 & C & M & +0,0000 & +0,0000 & 0 & +0,0008 & +0,0019 & 63 \\
Gemini-001 & C & R & +0,0000 & +0,0000 & 0 & +0,0000 & +0,0000 & 0 \\
\midrule
OpenAI-3-large & Q & ES & $-$0,0002 & +0,0003 & 2 & +0,0000 & +0,0000 & 0 \\
OpenAI-3-large & Q & M & +0,0000 & +0,0000 & 0 & $-$0,0008 & $-$0,0004 & 41 \\
OpenAI-3-large & Q & R & +0,0000 & +0,0000 & 0 & $-$0,0006 & $-$0,0017 & 40 \\
OpenAI-3-large & C & ES & +0,0002 & +0,0012 & 2 & $-$0,0011 & $-$0,0030 & 33 \\
OpenAI-3-large & C & M & +0,0000 & +0,0000 & 0 & +0,0007 & +0,0005 & 40 \\
OpenAI-3-large & C & R & +0,0000 & +0,0000 & 0 & +0,0000 & +0,0000 & 0 \\
\midrule
OpenAI-3-small & Q & ES & +0,0000 & +0,0000 & 0 & +0,0000 & +0,0000 & 0 \\
OpenAI-3-small & Q & M & $-$0,0002 & $-$0,0002 & 4 & $-$0,0035 & $-$0,0031 & 56 \\
OpenAI-3-small & Q & R & +0,0000 & +0,0000 & 0 & $-$0,0032 & $-$0,0022 & 56 \\
OpenAI-3-small & C & ES & +0,0000 & +0,0000 & 0 & +0,0000 & +0,0000 & 0 \\
OpenAI-3-small & C & M & +0,0000 & +0,0000 & 0 & +0,0000 & +0,0000 & 0 \\
OpenAI-3-small & C & R & +0,0000 & +0,0000 & 0 & $-$0,0017 & +0,0017 & 28 \\
\midrule
MiniLM-L6 & Q & ES & +0,0007 & +0,0006 & 9 & +0,0030 & +0,0007 & 37 \\
MiniLM-L6 & Q & M & +0,0000 & +0,0000 & 0 & +0,0012 & +0,0009 & 33 \\
MiniLM-L6 & Q & R & +0,0007 & +0,0017 & 6 & +0,0000 & +0,0000 & 0 \\
MiniLM-L6 & C & ES & +0,0000 & +0,0000 & 0 & +0,0043 & +0,0052 & 23 \\
MiniLM-L6 & C & M & +0,0000 & +0,0000 & 0 & +0,0034 & +0,0086 & 43 \\
MiniLM-L6 & C & R & +0,0000 & +0,0000 & 0 & +0,0030 & +0,0079 & 44 \\
\bottomrule
\end{tabular}
\caption{Ablation du \textit{gate} et de la \textit{search strictness}, NFCorpus. Variations calculées sur l'ensemble des requêtes du découpage initial. Var. : Q globale, C par clustering ; SS : \textit{search strictness} (ES \textit{error-scaled}, M \textit{mid}, R \textit{relaxed}) ; Dég. : requêtes initialement réussies devenues en erreur.}
\label{tab:s-ablation-nfcorpus}
\end{table}

\begin{table}[p]
\centering
\small
\begin{tabular}{lllrrrrrr}
\toprule
 & & & \multicolumn{3}{c}{\textit{Gate} strict} & \multicolumn{3}{c}{\textit{Gate} adaptatif} \\
\cmidrule(lr){4-6}\cmidrule(lr){7-9}
Encodeur & Var. & SS & $\Delta$NDCG@10 & $\Delta$MRR & Dég. & $\Delta$NDCG@10 & $\Delta$MRR & Dég. \\
\midrule
Titan-v1 & Q & ES & +0,0012 & +0,0011 & 0 & +0,0011 & +0,0022 & 0 \\
Titan-v1 & Q & M & $-$0,0001 & +0,0017 & 0 & $-$0,0003 & +0,0002 & 0 \\
Titan-v1 & Q & R & $-$0,0006 & +0,0005 & 0 & +0,0000 & +0,0000 & 0 \\
Titan-v1 & C & ES & +0,0023 & +0,0059 & 0 & +0,0009 & +0,0017 & 0 \\
Titan-v1 & C & M & +0,0035 & +0,0094 & 1 & +0,0001 & +0,0012 & 0 \\
Titan-v1 & C & R & +0,0002 & +0,0012 & 0 & +0,0001 & +0,0010 & 0 \\
\midrule
Gemini-001 & Q & ES & +0,0035 & +0,0040 & 6 & +0,0249 & +0,0336 & 21 \\
Gemini-001 & Q & M & +0,0066 & +0,0074 & 5 & +0,0220 & +0,0283 & 12 \\
Gemini-001 & Q & R & +0,0085 & +0,0114 & 7 & +0,0035 & +0,0046 & 4 \\
Gemini-001 & C & ES & +0,0084 & +0,0089 & 11 & +0,0313 & +0,0446 & 25 \\
Gemini-001 & C & M & +0,0169 & +0,0217 & 13 & +0,0021 & +0,0041 & 3 \\
Gemini-001 & C & R & +0,0000 & +0,0000 & 0 & +0,0154 & +0,0194 & 17 \\
\midrule
OpenAI-3-large & Q & ES & +0,0110 & +0,0195 & 3 & +0,0092 & +0,0135 & 3 \\
OpenAI-3-large & Q & M & +0,0056 & +0,0060 & 8 & +0,0012 & +0,0017 & 1 \\
OpenAI-3-large & Q & R & +0,0042 & +0,0046 & 6 & +0,0008 & +0,0023 & 1 \\
OpenAI-3-large & C & ES & +0,0068 & +0,0112 & 3 & $-$0,0001 & +0,0006 & 0 \\
OpenAI-3-large & C & M & +0,0121 & +0,0188 & 12 & +0,0000 & +0,0000 & 0 \\
OpenAI-3-large & C & R & +0,0031 & +0,0051 & 6 & +0,0015 & +0,0041 & 0 \\
\midrule
OpenAI-3-small & Q & ES & +0,0051 & +0,0113 & 4 & +0,0038 & +0,0061 & 5 \\
OpenAI-3-small & Q & M & +0,0031 & +0,0073 & 6 & +0,0027 & +0,0051 & 5 \\
OpenAI-3-small & Q & R & +0,0029 & +0,0084 & 6 & +0,0018 & +0,0056 & 2 \\
OpenAI-3-small & C & ES & +0,0063 & +0,0113 & 7 & +0,0000 & +0,0000 & 0 \\
OpenAI-3-small & C & M & +0,0045 & +0,0077 & 5 & +0,0000 & +0,0010 & 0 \\
OpenAI-3-small & C & R & +0,0042 & +0,0103 & 10 & +0,0000 & +0,0000 & 0 \\
\midrule
MiniLM-L6 & Q & ES & +0,0000 & +0,0000 & 0 & +0,0000 & +0,0000 & 0 \\
MiniLM-L6 & Q & M & $-$0,0001 & +0,0011 & 2 & +0,0000 & +0,0000 & 0 \\
MiniLM-L6 & Q & R & $-$0,0014 & $-$0,0004 & 6 & $-$0,0004 & +0,0028 & 4 \\
MiniLM-L6 & C & ES & +0,0000 & +0,0000 & 0 & +0,0000 & +0,0000 & 0 \\
MiniLM-L6 & C & M & $-$0,0016 & $-$0,0009 & 8 & +0,0000 & +0,0000 & 0 \\
MiniLM-L6 & C & R & $-$0,0005 & +0,0006 & 6 & +0,0000 & +0,0000 & 0 \\
\bottomrule
\end{tabular}
\caption{Ablation du \textit{gate} et de la \textit{search strictness}, SciDocs. Variations calculées sur l'ensemble des requêtes du découpage initial. Var. : Q globale, C par clustering ; SS : \textit{search strictness} (ES \textit{error-scaled}, M \textit{mid}, R \textit{relaxed}) ; Dég. : requêtes initialement réussies devenues en erreur.}
\label{tab:s-ablation-scidocs}
\end{table}

\begin{table}[p]
\centering
\small
\begin{tabular}{lllrrrrrr}
\toprule
 & & & \multicolumn{3}{c}{\textit{Gate} strict} & \multicolumn{3}{c}{\textit{Gate} adaptatif} \\
\cmidrule(lr){4-6}\cmidrule(lr){7-9}
Encodeur & Var. & SS & $\Delta$NDCG@10 & $\Delta$MRR & Dég. & $\Delta$NDCG@10 & $\Delta$MRR & Dég. \\
\midrule
Titan-v1 & Q & ES & +0,0008 & +0,0011 & 2 & +0,0000 & +0,0000 & 0 \\
Titan-v1 & Q & M & +0,0054 & +0,0043 & 2 & +0,0000 & +0,0000 & 0 \\
Titan-v1 & Q & R & +0,0116 & +0,0095 & 0 & +0,0000 & +0,0000 & 0 \\
Titan-v1 & C & ES & +0,0100 & +0,0108 & 0 & +0,0000 & +0,0000 & 0 \\
Titan-v1 & C & M & +0,0054 & +0,0056 & 0 & +0,0000 & +0,0000 & 0 \\
Titan-v1 & C & R & +0,0137 & +0,0101 & 0 & +0,0000 & +0,0000 & 0 \\
\midrule
Gemini-001 & Q & ES & +0,0010 & +0,0021 & 1 & +0,0000 & +0,0000 & 0 \\
Gemini-001 & Q & M & +0,0010 & +0,0021 & 1 & +0,0000 & +0,0000 & 0 \\
Gemini-001 & Q & R & +0,0016 & +0,0029 & 2 & +0,0000 & +0,0000 & 0 \\
Gemini-001 & C & ES & +0,0019 & +0,0029 & 1 & +0,0000 & +0,0000 & 0 \\
Gemini-001 & C & M & +0,0014 & +0,0019 & 0 & +0,0000 & +0,0000 & 0 \\
Gemini-001 & C & R & +0,0019 & +0,0029 & 1 & +0,0000 & +0,0000 & 0 \\
\midrule
OpenAI-3-large & Q & ES & +0,0214 & +0,0239 & 3 & +0,0067 & +0,0050 & 4 \\
OpenAI-3-large & Q & M & +0,0181 & +0,0180 & 5 & +0,0094 & +0,0080 & 5 \\
OpenAI-3-large & Q & R & +0,0074 & +0,0056 & 7 & +0,0084 & +0,0069 & 6 \\
OpenAI-3-large & C & ES & +0,0469 & +0,0450 & 2 & +0,0103 & +0,0092 & 4 \\
OpenAI-3-large & C & M & +0,0119 & +0,0094 & 6 & +0,0032 & +0,0029 & 6 \\
OpenAI-3-large & C & R & +0,0188 & +0,0162 & 5 & +0,0042 & +0,0040 & 7 \\
\midrule
OpenAI-3-small & Q & ES & +0,0035 & +0,0056 & 1 & +0,0000 & +0,0000 & 0 \\
OpenAI-3-small & Q & M & $-$0,0014 & $-$0,0005 & 5 & +0,0000 & +0,0000 & 0 \\
OpenAI-3-small & Q & R & $-$0,0001 & +0,0008 & 3 & +0,0000 & +0,0000 & 0 \\
OpenAI-3-small & C & ES & +0,0009 & +0,0010 & 4 & +0,0000 & +0,0000 & 0 \\
OpenAI-3-small & C & M & +0,0072 & +0,0069 & 5 & +0,0000 & +0,0000 & 0 \\
OpenAI-3-small & C & R & +0,0042 & +0,0030 & 8 & +0,0000 & +0,0000 & 0 \\
\midrule
MiniLM-L6 & Q & ES & +0,0115 & +0,0139 & 1 & +0,0000 & +0,0000 & 0 \\
MiniLM-L6 & Q & M & +0,0071 & +0,0068 & 4 & +0,0000 & +0,0000 & 0 \\
MiniLM-L6 & Q & R & +0,0048 & +0,0039 & 5 & +0,0000 & +0,0000 & 0 \\
MiniLM-L6 & C & ES & +0,0107 & +0,0116 & 3 & +0,0000 & +0,0000 & 0 \\
MiniLM-L6 & C & M & +0,0182 & +0,0204 & 4 & +0,0000 & +0,0000 & 0 \\
MiniLM-L6 & C & R & +0,0199 & +0,0199 & 7 & +0,0000 & +0,0000 & 0 \\
\bottomrule
\end{tabular}
\caption{Ablation du \textit{gate} et de la \textit{search strictness}, SciFact. Variations calculées sur l'ensemble des requêtes du découpage initial. Var. : Q globale, C par clustering ; SS : \textit{search strictness} (ES \textit{error-scaled}, M \textit{mid}, R \textit{relaxed}) ; Dég. : requêtes initialement réussies devenues en erreur.}
\label{tab:s-ablation-scifact}
\end{table}

\begin{table}[p]
\centering
\small
\begin{tabular}{lllrrr}
\toprule
Encodeur & \textit{Gate} & SS & $k=3$ & $k=5$ & $k=10$ \\
\midrule
Titan-v1 & S & ES & +0,0415 & +0,0400 & +0,0354 \\
Titan-v1 & S & M & +0,0508 & +0,0392 & +0,0364 \\
Titan-v1 & S & R & +0,0522 & +0,0387 & +0,0328 \\
Titan-v1 & A & ES & +0,0000 & +0,0000 & +0,0000 \\
Titan-v1 & A & M & +0,0000 & +0,0000 & +0,0000 \\
Titan-v1 & A & R & +0,0000 & +0,0000 & +0,0000 \\
\midrule
Gemini-001 & S & ES & +0,0000 & +0,0011 & +0,0225 \\
Gemini-001 & S & M & +0,0480 & +0,0367 & +0,0285 \\
Gemini-001 & S & R & +0,0370 & +0,0525 & +0,0312 \\
Gemini-001 & A & ES & +0,1459 & +0,1330 & +0,1508 \\
Gemini-001 & A & M & +0,0881 & +0,0778 & +0,0976 \\
Gemini-001 & A & R & +0,0977 & +0,1149 & +0,0954 \\
\midrule
OpenAI-3-large & S & ES & +0,0788 & +0,0819 & +0,0736 \\
OpenAI-3-large & S & M & +0,0870 & +0,0872 & +0,0771 \\
OpenAI-3-large & S & R & +0,0870 & +0,0857 & +0,0801 \\
OpenAI-3-large & A & ES & +0,0000 & +0,0000 & +0,0000 \\
OpenAI-3-large & A & M & +0,0000 & +0,0000 & +0,0000 \\
OpenAI-3-large & A & R & +0,0000 & +0,0000 & +0,0000 \\
\midrule
OpenAI-3-small & S & ES & +0,0814 & +0,0763 & +0,0658 \\
OpenAI-3-small & S & M & +0,0784 & +0,0806 & +0,0616 \\
OpenAI-3-small & S & R & +0,0820 & +0,0751 & +0,0633 \\
OpenAI-3-small & A & ES & +0,0000 & +0,0000 & +0,0000 \\
OpenAI-3-small & A & M & +0,0000 & +0,0000 & +0,0000 \\
OpenAI-3-small & A & R & +0,0000 & +0,0000 & +0,0000 \\
\midrule
MiniLM-L6 & S & ES & +0,0500 & +0,0535 & +0,0536 \\
MiniLM-L6 & S & M & +0,0496 & +0,0522 & +0,0482 \\
MiniLM-L6 & S & R & +0,0526 & +0,0592 & +0,0465 \\
MiniLM-L6 & A & ES & +0,0000 & +0,0000 & +0,0000 \\
MiniLM-L6 & A & M & +0,0000 & +0,0000 & +0,0000 \\
MiniLM-L6 & A & R & +0,0000 & +0,0000 & +0,0000 \\
\bottomrule
\end{tabular}
\caption{Sensibilité au nombre de clusters $k$ de la variante par clustering, ArguAna. Variations de NDCG@10 calculées sur l'ensemble des requêtes du découpage initial. \textit{Gate} : S strict, A adaptatif ; SS : \textit{search strictness} (ES, M, R).}
\label{tab:s-clusters-arguana}
\end{table}

\begin{table}[p]
\centering
\small
\begin{tabular}{lllrrr}
\toprule
Encodeur & \textit{Gate} & SS & $k=3$ & $k=5$ & $k=10$ \\
\midrule
Titan-v1 & S & ES & +0,0000 & +0,0000 & +0,0000 \\
Titan-v1 & S & M & +0,0012 & +0,0018 & +0,0000 \\
Titan-v1 & S & R & +0,0000 & +0,0012 & +0,0007 \\
Titan-v1 & A & ES & +0,0043 & +0,0005 & +0,0002 \\
Titan-v1 & A & M & +0,0001 & +0,0001 & +0,0001 \\
Titan-v1 & A & R & +0,0000 & +0,0003 & +0,0000 \\
\midrule
Gemini-001 & S & ES & +0,0000 & +0,0000 & +0,0000 \\
Gemini-001 & S & M & +0,0000 & +0,0000 & +0,0000 \\
Gemini-001 & S & R & +0,0000 & +0,0000 & +0,0006 \\
Gemini-001 & A & ES & +0,0041 & +0,0031 & +0,0039 \\
Gemini-001 & A & M & +0,0008 & +0,0030 & +0,0017 \\
Gemini-001 & A & R & +0,0036 & +0,0038 & +0,0000 \\
\midrule
OpenAI-3-large & S & ES & $-$0,0002 & +0,0002 & $-$0,0001 \\
OpenAI-3-large & S & M & +0,0000 & +0,0001 & +0,0000 \\
OpenAI-3-large & S & R & +0,0000 & +0,0000 & +0,0000 \\
OpenAI-3-large & A & ES & $-$0,0011 & +0,0032 & $-$0,0002 \\
OpenAI-3-large & A & M & +0,0006 & $-$0,0008 & +0,0007 \\
OpenAI-3-large & A & R & +0,0000 & +0,0000 & +0,0001 \\
\midrule
OpenAI-3-small & S & ES & +0,0009 & +0,0006 & +0,0000 \\
OpenAI-3-small & S & M & +0,0005 & +0,0003 & +0,0000 \\
OpenAI-3-small & S & R & +0,0000 & $-$0,0001 & +0,0000 \\
OpenAI-3-small & A & ES & +0,0006 & +0,0007 & +0,0000 \\
OpenAI-3-small & A & M & +0,0000 & +0,0002 & $-$0,0008 \\
OpenAI-3-small & A & R & +0,0008 & +0,0002 & $-$0,0017 \\
\midrule
MiniLM-L6 & S & ES & +0,0000 & +0,0000 & +0,0000 \\
MiniLM-L6 & S & M & +0,0004 & +0,0002 & +0,0000 \\
MiniLM-L6 & S & R & +0,0007 & +0,0000 & +0,0004 \\
MiniLM-L6 & A & ES & +0,0043 & +0,0029 & +0,0027 \\
MiniLM-L6 & A & M & +0,0034 & +0,0003 & +0,0005 \\
MiniLM-L6 & A & R & +0,0030 & +0,0004 & +0,0005 \\
\bottomrule
\end{tabular}
\caption{Sensibilité au nombre de clusters $k$ de la variante par clustering, NFCorpus. Variations de NDCG@10 calculées sur l'ensemble des requêtes du découpage initial. \textit{Gate} : S strict, A adaptatif ; SS : \textit{search strictness} (ES, M, R).}
\label{tab:s-clusters-nfcorpus}
\end{table}

\begin{table}[p]
\centering
\small
\begin{tabular}{lllrrr}
\toprule
Encodeur & \textit{Gate} & SS & $k=3$ & $k=5$ & $k=10$ \\
\midrule
Titan-v1 & S & ES & +0,0013 & +0,0001 & +0,0023 \\
Titan-v1 & S & M & +0,0035 & +0,0017 & +0,0015 \\
Titan-v1 & S & R & +0,0005 & +0,0002 & +0,0002 \\
Titan-v1 & A & ES & +0,0009 & +0,0010 & +0,0005 \\
Titan-v1 & A & M & +0,0001 & +0,0000 & +0,0000 \\
Titan-v1 & A & R & +0,0000 & +0,0000 & +0,0001 \\
\midrule
Gemini-001 & S & ES & +0,0084 & +0,0081 & +0,0105 \\
Gemini-001 & S & M & +0,0169 & +0,0028 & +0,0043 \\
Gemini-001 & S & R & +0,0050 & +0,0000 & +0,0169 \\
Gemini-001 & A & ES & +0,0249 & +0,0313 & +0,0140 \\
Gemini-001 & A & M & +0,0021 & +0,0006 & +0,0024 \\
Gemini-001 & A & R & +0,0154 & $-$0,0002 & +0,0139 \\
\midrule
OpenAI-3-large & S & ES & +0,0068 & +0,0087 & +0,0033 \\
OpenAI-3-large & S & M & +0,0087 & +0,0002 & +0,0121 \\
OpenAI-3-large & S & R & +0,0012 & +0,0031 & +0,0126 \\
OpenAI-3-large & A & ES & +0,0000 & +0,0023 & $-$0,0001 \\
OpenAI-3-large & A & M & +0,0001 & +0,0000 & +0,0000 \\
OpenAI-3-large & A & R & +0,0015 & +0,0013 & +0,0022 \\
\midrule
OpenAI-3-small & S & ES & +0,0063 & +0,0041 & +0,0067 \\
OpenAI-3-small & S & M & +0,0020 & +0,0045 & +0,0037 \\
OpenAI-3-small & S & R & +0,0042 & +0,0028 & +0,0023 \\
OpenAI-3-small & A & ES & +0,0058 & +0,0033 & +0,0000 \\
OpenAI-3-small & A & M & +0,0001 & +0,0010 & +0,0000 \\
OpenAI-3-small & A & R & +0,0000 & +0,0001 & +0,0040 \\
\midrule
MiniLM-L6 & S & ES & +0,0007 & $-$0,0003 & +0,0000 \\
MiniLM-L6 & S & M & +0,0001 & $-$0,0016 & $-$0,0017 \\
MiniLM-L6 & S & R & $-$0,0005 & +0,0008 & $-$0,0023 \\
MiniLM-L6 & A & ES & +0,0000 & $-$0,0005 & +0,0000 \\
MiniLM-L6 & A & M & +0,0000 & +0,0000 & $-$0,0003 \\
MiniLM-L6 & A & R & $-$0,0008 & $-$0,0012 & +0,0000 \\
\bottomrule
\end{tabular}
\caption{Sensibilité au nombre de clusters $k$ de la variante par clustering, SciDocs. Variations de NDCG@10 calculées sur l'ensemble des requêtes du découpage initial. \textit{Gate} : S strict, A adaptatif ; SS : \textit{search strictness} (ES, M, R).}
\label{tab:s-clusters-scidocs}
\end{table}

\begin{table}[p]
\centering
\small
\begin{tabular}{lllrrr}
\toprule
Encodeur & \textit{Gate} & SS & $k=3$ & $k=5$ & $k=10$ \\
\midrule
Titan-v1 & S & ES & +0,0054 & +0,0182 & +0,0100 \\
Titan-v1 & S & M & +0,0306 & +0,0032 & +0,0054 \\
Titan-v1 & S & R & +0,0722 & +0,0137 & +0,0028 \\
Titan-v1 & A & ES & +0,0000 & +0,0000 & +0,0000 \\
Titan-v1 & A & M & +0,0000 & +0,0000 & +0,0000 \\
Titan-v1 & A & R & +0,0000 & +0,0000 & +0,0000 \\
\midrule
Gemini-001 & S & ES & +0,0019 & -- & -- \\
Gemini-001 & S & M & +0,0014 & -- & -- \\
Gemini-001 & S & R & +0,0019 & -- & -- \\
Gemini-001 & A & ES & +0,0000 & -- & -- \\
Gemini-001 & A & M & +0,0000 & -- & -- \\
Gemini-001 & A & R & +0,0000 & -- & -- \\
\midrule
OpenAI-3-large & S & ES & +0,0385 & +0,0469 & +0,0518 \\
OpenAI-3-large & S & M & +0,0119 & +0,0067 & +0,0290 \\
OpenAI-3-large & S & R & +0,0188 & +0,0061 & +0,0358 \\
OpenAI-3-large & A & ES & +0,0103 & +0,0001 & +0,0226 \\
OpenAI-3-large & A & M & +0,0099 & +0,0032 & +0,0054 \\
OpenAI-3-large & A & R & +0,0056 & +0,0042 & +0,0133 \\
\midrule
OpenAI-3-small & S & ES & +0,0007 & +0,0009 & +0,0007 \\
OpenAI-3-small & S & M & +0,0031 & +0,0012 & +0,0072 \\
OpenAI-3-small & S & R & +0,0042 & +0,0025 & +0,0084 \\
OpenAI-3-small & A & ES & +0,0000 & +0,0000 & +0,0000 \\
OpenAI-3-small & A & M & +0,0000 & +0,0000 & +0,0000 \\
OpenAI-3-small & A & R & +0,0000 & +0,0000 & +0,0000 \\
\midrule
MiniLM-L6 & S & ES & +0,0014 & +0,0107 & +0,0111 \\
MiniLM-L6 & S & M & +0,0035 & +0,0196 & +0,0182 \\
MiniLM-L6 & S & R & +0,0051 & +0,0121 & +0,0199 \\
MiniLM-L6 & A & ES & +0,0000 & +0,0000 & +0,0000 \\
MiniLM-L6 & A & M & +0,0000 & +0,0000 & +0,0000 \\
MiniLM-L6 & A & R & +0,0000 & +0,0000 & +0,0000 \\
\bottomrule
\end{tabular}
\caption{Sensibilité au nombre de clusters $k$ de la variante par clustering, SciFact. Variations de NDCG@10 calculées sur l'ensemble des requêtes du découpage initial. \textit{Gate} : S strict, A adaptatif ; SS : \textit{search strictness} (ES, M, R). -- : configuration non exécutable, le nombre d'erreurs résiduelles étant inférieur au nombre de clusters demandé.}
\label{tab:s-clusters-scifact}
\end{table}


\section{Coût de calcul et environnement}

Les exécutions ont été réalisées sur une machine virtuelle sans GPU (64 vCPU Intel Xeon Ice Lake à 2,9\,GHz, 256\,Gio de RAM, Ubuntu Linux). Chaque exécution utilisait un seul fil de calcul BLAS (variables \texttt{OMP\_NUM\_THREADS} et \texttt{MKL\_NUM\_THREADS} fixées à~1) ; les exécutions indépendantes étaient lancées en parallèle. L'implémentation, en Python~3, utilise NumPy, PyTorch (version CPU), scikit-learn et \texttt{pytrec\_eval} ; la liste des dépendances est fournie avec le code, sans versions verrouillées. Les découpages, le clustering des erreurs et l'initialisation des directions utilisent une graine fixe : 42 pour le découpage initial, 43 à 46 pour les découpages supplémentaires.

\begin{table}[p]
\centering
\small
\begin{tabular}{llrrrr}
\toprule
Jeu & Encodeur & Exéc. globales & Durée moy. & Exéc. clustering & Durée moy. \\
\midrule
ArguAna & Titan-v1 & 6 & 10,6 & 18 & 28,7 \\
ArguAna & Gemini-001 & 6 & 10,5 & 18 & 13,8 \\
ArguAna & OpenAI-3-large & 6 & 6,7 & 18 & 42,0 \\
ArguAna & OpenAI-3-small & 6 & 6,1 & 18 & 37,3 \\
ArguAna & MiniLM-L6 & 6 & 8,5 & 18 & 34,0 \\
\midrule
NFCorpus & Titan-v1 & 6 & 41,2 & 18 & 18,1 \\
NFCorpus & Gemini-001 & 6 & 23,4 & 18 & 9,7 \\
NFCorpus & OpenAI-3-large & 6 & 23,1 & 18 & 11,2 \\
NFCorpus & OpenAI-3-small & 6 & 25,5 & 18 & 11,4 \\
NFCorpus & MiniLM-L6 & 6 & 32,9 & 18 & 12,0 \\
\midrule
SciDocs & Titan-v1 & 6 & 32,0 & 18 & 20,0 \\
SciDocs & Gemini-001 & 6 & 30,4 & 18 & 15,8 \\
SciDocs & OpenAI-3-large & 6 & 29,3 & 18 & 20,3 \\
SciDocs & OpenAI-3-small & 6 & 32,1 & 18 & 20,5 \\
SciDocs & MiniLM-L6 & 6 & 19,5 & 18 & 15,5 \\
\midrule
SciFact & Titan-v1 & 6 & 5,6 & 18 & 9,5 \\
SciFact & Gemini-001 & 6 & 15,6 & 6 & 2,8 \\
SciFact & OpenAI-3-large & 6 & 5,2 & 18 & 17,5 \\
SciFact & OpenAI-3-small & 6 & 3,5 & 18 & 7,4 \\
SciFact & MiniLM-L6 & 6 & 5,7 & 18 & 7,0 \\
\midrule
Tous les couples &  & 120 & 18,4 & 348 & 18,2 \\
\bottomrule
\end{tabular}
\caption{Coût de calcul : durée moyenne par exécution aboutie, en minutes, selon la variante. Les exécutions étant indépendantes et parallélisées, la somme des durées ne correspond pas à la durée totale de la campagne.}
\label{tab:s-durees}
\end{table}


\section{Dimension des embeddings, taille du corpus et taux d'erreur}

Le tableau~\ref{tab:s-covariables} rapporte, pour la configuration retenue sur le découpage initial, le gain de NDCG@10 en test sur ce découpage. En moyenne par dimension, il vaut $+0{,}004$ à 384 dimensions (4 couples), $+0{,}013$ à 1\,024 (4 couples), $+0{,}014$ à 1\,536 (4 couples) et $+0{,}035$ à 3\,072 (8 couples). Sur les vingt couples, la corrélation de Spearman avec le gain vaut $+0{,}34$ pour le taux d'erreur initial et $+0{,}23$ pour la dimension (Pearson $+0{,}34$). La taille du corpus ne montre pas de relation cohérente (Pearson $-0{,}09$, Spearman $+0{,}28$).

Ces observations appellent trois réserves. D'abord, la dimension est confondue avec la famille de modèles : le groupe à 3\,072 dimensions réunit précisément les deux modèles commerciaux les plus performants (Gemini et OpenAI-3-large). Ensuite, le taux d'erreur initial est lui-même lié à la qualité du modèle : SciFact--Gemini, avec 2\,\% d'erreurs initiales, ne laisse presque aucune marge de correction malgré ses 3\,072 dimensions. Enfin, avec vingt couples et un seul découpage, ces corrélations décrivent les observations sans permettre de conclusion statistique par variable.

\begin{table}[p]
\centering
\small
\begin{tabular}{llrrrrrr}
\toprule
Jeu & Encodeur & Dim. & Corpus & NDCG@10 & MRR & Erreurs & $\Delta$NDCG@10 test \\
\midrule
ArguAna & Gemini-001 & 3072 & 8674 & 0,5762 & 0,4565 & 0,89 & +0,1626 \\
ArguAna & OpenAI-3-large & 3072 & 8674 & 0,4100 & 0,2846 & 1,00 & +0,0535 \\
ArguAna & OpenAI-3-small & 1536 & 8674 & 0,4018 & 0,2802 & 1,00 & +0,0506 \\
SciFact & Titan-v1 & 1024 & 5183 & 0,7324 & 0,7019 & 0,39 & +0,0323 \\
SciFact & OpenAI-3-large & 3072 & 5183 & 0,7700 & 0,7409 & 0,35 & +0,0278 \\
SciDocs & Gemini-001 & 3072 & 25657 & 0,2411 & 0,4025 & 0,73 & +0,0220 \\
ArguAna & Titan-v1 & 1024 & 8674 & 0,3654 & 0,2567 & 1,00 & +0,0116 \\
ArguAna & MiniLM-L6 & 384 & 8674 & 0,3727 & 0,2592 & 1,00 & +0,0113 \\
SciDocs & OpenAI-3-large & 3072 & 25657 & 0,2231 & 0,3767 & 0,75 & +0,0099 \\
SciFact & MiniLM-L6 & 384 & 5183 & 0,6488 & 0,6182 & 0,49 & +0,0077 \\
SciFact & OpenAI-3-small & 1536 & 5183 & 0,7221 & 0,6868 & 0,41 & +0,0065 \\
NFCorpus & Titan-v1 & 1024 & 3633 & 0,3257 & 0,5441 & 0,54 & +0,0043 \\
SciDocs & Titan-v1 & 1024 & 25657 & 0,2072 & 0,3591 & 0,77 & +0,0042 \\
NFCorpus & Gemini-001 & 3072 & 3633 & 0,4097 & 0,6235 & 0,47 & +0,0041 \\
SciDocs & OpenAI-3-small & 1536 & 25657 & 0,1922 & 0,3367 & 0,78 & +0,0018 \\
NFCorpus & MiniLM-L6 & 384 & 3633 & 0,3163 & 0,5148 & 0,58 & +0,0017 \\
NFCorpus & OpenAI-3-large & 3072 & 3633 & 0,3975 & 0,6123 & 0,48 & +0,0012 \\
SciFact & Gemini-001 & 3072 & 5183 & 0,9931 & 0,9903 & 0,02 & +0,0000 \\
NFCorpus & OpenAI-3-small & 1536 & 3633 & 0,3598 & 0,5668 & 0,54 & $-$0,0037 \\
SciDocs & MiniLM-L6 & 384 & 25657 & 0,2040 & 0,3542 & 0,78 & $-$0,0040 \\
\bottomrule
\end{tabular}
\caption{Caractéristiques des couples pour la configuration retenue sur le découpage initial, triées par gain de NDCG@10 en test sur ce découpage. Dim. : dimension des embeddings ; Corpus : nombre de documents ; NDCG@10 et MRR : valeurs avant correction ; Erreurs : part initiale des requêtes dont le premier résultat n'est pas pertinent.}
\label{tab:s-covariables}
\end{table}


\end{document}
